\section{Related Work}

The goal of multimodal fake news detection is to identify inconsistencies or deceptive content across modalities such as text and images. BMR~\cite{ying2023bootstrapping} employ single-view prediction to disentangle multimodal features and explain which modality plays a critical role in the final decision. NSLM~\cite{dong2024unveiling} leverages symbolic logic constraints to unveil deceptive patterns within multimodal fake news. MMDFND~\cite{tong2024mmdfnd} enhances multimodal feature representation by integrating additional vision and language models and applying feature-wise weighting strategies. FKA-Owl~\cite{liu2024fka} encodes forgery-specific features as learnable continuous vectors and utilizes soft prompt tuning to strengthen the reasoning ability of LVLMs. Although these approaches improve the effectiveness of multimodal fake news detection, they remain limited by challenges in modality alignment and a strong reliance on high-quality deception-specific features. Retrieval-Augmented Generation (RAG) alleviates these issues by incorporating relevant information retrieved from external corpora~\cite{rag1,rag2,rag3,rag4,rag5,rag6}. More recently, dynamic RAG techniques further enhance this capability by adaptively deciding when to retrieve external evidence and what to retrieve during inference, thereby improving model performance and task accuracy~\cite{drag1,drag3,drag4,drag5}. Building upon this line of work, we introduce an improved dynamic RAG framework that enables accurate and efficient identification of deceptive cues, leading to more effective and efficient multimodal fake news detection.
